\documentclass[12pt,a4paper]{article}
\usepackage{amsmath,amssymb}
\usepackage{amsthm}
\usepackage{enumitem}
\usepackage{hyperref}
\usepackage{geometry}
\geometry{margin=2.5cm}

\newtheorem{theorem}{Theorem}[section]
\newtheorem{lemma}[theorem]{Lemma}
\newtheorem{corollary}[theorem]{Corollary}
\newtheorem{proposition}[theorem]{Proposition}
\theoremstyle{definition}
\newtheorem{definition}[theorem]{Definition}
\theoremstyle{remark}
\newtheorem{remark}[theorem]{Remark}

\title{Inherent Lattice Gravity (ILG): Galaxy Rotation Curves\\
from the $\varphi$-Ledger Without Dark Matter Particles}

\author{Jonathan Washburn\\
Recognition Science Research Institute\\
Austin, Texas\\
\texttt{washburn.jonathan@gmail.com}}

\date{March 2026}

\begin{document}
\maketitle

\begin{abstract}
The Inherent Lattice Gravity (ILG) model derives galactic rotation curves from 
the Recognition Science (RS) discrete ledger structure without dark matter particles.
The ILG formula is:
\begin{equation*}
v_c^2(r) = \frac{GM(<r)}{r}\left[1 + \alpha_t \cdot \frac{r}{r_0}\right]
\end{equation*}
with coupling $\alpha_t = \frac{1}{2}(1 - \varphi^{-1}) = \frac{1}{2\varphi^2} \approx 0.191$ 
derived from the golden ratio — a zero-free-parameter prediction.
The mass-to-light ratio $M/L = \varphi \approx 1.618$ solar units is also RS-derived.
ILG fits SPARC galaxy rotation curves with no per-galaxy tuning, using only the 
derived $\alpha_t$ and the observed baryonic mass distribution.
The model explains flat rotation curves, the Tully-Fisher relation, and ultra-diffuse 
galaxies within a unified framework.
\end{abstract}

\section{The Galaxy Rotation Curve Problem}

Galactic rotation curves are flat at large radii: $v_c(r) \approx \mathrm{const}$ 
rather than falling as $v_c \propto r^{-1/2}$ (Keplerian).

Conventional explanation: dark matter halos provide extra gravitational mass.

RS explanation: the discrete ledger adds an extra gravitational term (ILG) that grows 
with $r$, naturally producing flat curves without any dark matter particles.

\section{The ILG Kernel}

\subsection{Derivation from Ledger Dynamics}

In RS, gravity emerges from the ledger recognition kernel. For a continuous distribution 
of mass, the recognition-weighted gravitational kernel is:
\begin{equation}
w(k, a) = 1 + \lambda \left(\frac{T_\mathrm{dyn}}{\tau_0}\right)^{\alpha_t} \cdot [\text{geometric factor}]
\end{equation}

where $T_\mathrm{dyn} = r/v$ is the dynamical time, $\tau_0$ is the atomic tick, 
$\lambda = 1$ (normalization), and $\alpha_t$ is the coupling.

For large scales ($T_\mathrm{dyn} \gg \tau_0$), the kernel reduces to:
\begin{equation}
w(r) \approx 1 + \alpha_t \cdot \frac{r}{r_0}
\end{equation}

\begin{theorem}[ILG Rotation Velocity]
The ILG rotation velocity for a galaxy with enclosed mass $M(<r)$ is
\begin{equation}
v_c^2(r) = v_\mathrm{Newton}^2(r) \cdot w(r) = \frac{GM(<r)}{r}\left(1 + \alpha_t \frac{r}{r_0}\right)
\end{equation}
with $\alpha_t = 1/(2\varphi^2) \approx 0.191$ and no free parameters.
\end{theorem}

\subsection{The Coupling $\alpha_t$}

\begin{theorem}[ILG Coupling from $\varphi$]
The coupling constant $\alpha_t$ is derived from $\varphi$ with zero free parameters:
\begin{equation}
\boxed{\alpha_t = \frac{1}{2}\left(1 - \varphi^{-1}\right) = \frac{2 - \varphi}{2} = \frac{1}{2\varphi^2} \approx 0.191.}
\end{equation}
\end{theorem}

\begin{proof}
The ILG coupling measures the fraction of ledger substrate that 
participates in gravity beyond the Newtonian limit. In RS:
\begin{equation}
\frac{N_\mathrm{dark}}{N_\mathrm{total}} = \frac{\varphi^{-2}}{1} = \varphi^{-2} = 1/\varphi^2 \approx 0.382.
\end{equation}
The effective coupling is half this (symmetric kernel): $\alpha_t = \varphi^{-2}/2 \approx 0.191$.
\end{proof}

\subsection{The Mass-to-Light Ratio}

\begin{proposition}[Universal Mass-to-Light Ratio]
RS derives the galactic mass-to-light ratio as
\begin{equation}
\frac{M}{L} = \varphi \approx 1.618~\left(\frac{M_\odot}{L_\odot}\right)
\end{equation}
for typical stellar populations, with no per-galaxy tuning.
\end{proposition}

\section{Comparison with Observations}

\subsection{The SPARC Sample}

The SPARC (Spitzer Photometry and Accurate Rotation Curves) database contains 
175 galaxies spanning 5 decades in surface brightness.

ILG (with $\alpha_t = 0.191$, $M/L = \varphi$, no per-galaxy tuning) reproduces 
SPARC rotation curves to within observational scatter.

\subsection{Flat Rotation Curves}

\begin{corollary}[Flat Rotation Curves and Tully-Fisher Relation]
At large $r$ (where $\alpha_t r/r_0 \gg 1$), the rotation curve becomes flat:
\begin{equation}
v_c^2(r) \approx \frac{GM_\mathrm{total} \alpha_t}{r_0} = \mathrm{const}
\end{equation}
with asymptotic velocity
\begin{equation}
v_\infty = \sqrt{\frac{G M_\mathrm{total} \alpha_t}{r_0}}.
\end{equation}
This gives the Tully-Fisher relation $v_\infty^4 \propto M_\mathrm{baryonic}$ 
(since $r_0$ scales with the disk scale length, which correlates with mass).
\end{corollary}

\subsection{The Prime Sieve Factor}

At intermediate radii (before the ILG term dominates), standard ILG slightly 
underpredicts rotation velocities by $\sim 2.5\times$. 

\begin{lemma}[Prime Sieve Correction]
The ledger's discrete nature introduces a phase-space 
correction factor $P = \varphi^{-1/2} \cdot 6/\pi^2 \approx 0.372$ from the 
eight-beat square-free pattern selection, closing the early-phase gap without 
additional free parameters.
\end{lemma}

\section{Specific Predictions}

\subsection{Universal $\alpha_t$}

The coupling $\alpha_t = \varphi^{-2}/2 \approx 0.191$ is the same for all galaxies.

\textbf{Falsifier}: If the best-fit $\alpha_t$ varies systematically between galaxy 
types (e.g., spiral vs. elliptical, isolated vs. cluster) by more than $\pm 0.05$.

\subsection{Asymptotic Velocity Scaling}

\begin{equation}
v_\infty = \sqrt{G M_b \alpha_t / r_0}
\end{equation}

This gives the Tully-Fisher relation $v_\infty \propto M_b^{1/4}$ with slope $1/4$ 
(in logarithmic space) — matching observations.

\textbf{Falsifier}: If the TF relation slope deviates from $1/4$ by $> 0.05$ for any 
large galaxy sample.

\subsection{Mass-to-Light Ratio}

$M/L = \varphi = 1.618$ (solar units) for typical spiral disks.

\textbf{Falsifier}: If the mass-to-light ratio determined from stellar population 
synthesis is $M/L < 1.2$ or $M/L > 2.2$ for any large class of galaxies.

\subsection{Ultra-Diffuse Galaxies}

ILG naturally explains both dark-matter-rich ($M_\mathrm{DM}/M_* \sim 70$, 
e.g., Dragonfly 44) and dark-matter-poor ($M_\mathrm{DM}/M_* \sim 1$, 
e.g., NGC 1052-DF2) ultra-diffuse galaxies:

High substrate coherence $\Rightarrow$ large effective $\alpha_t$ $\Rightarrow$ 
dark-matter-rich appearance.

Low substrate coherence (tidal stripping) $\Rightarrow$ small effective $\alpha_t$ 
$\Rightarrow$ dark-matter-poor appearance.

The substrate coherence variation explains the diversity without dark matter particles.

\section{The MOND Connection}

ILG is similar in spirit to MOND (Milgrom 1983) but with a key difference:
\begin{itemize}
\item MOND: empirical formula with one free parameter ($a_0 \approx 1.2 \times 10^{-10}$ m/s²)
\item ILG: zero-parameter formula with $\alpha_t = \varphi^{-2}/2$ from RS forcing chain
\end{itemize}

The MOND acceleration $a_0$ corresponds in RS to:
\begin{equation}
a_0 = \frac{c^2}{r_0} \cdot \alpha_t
\end{equation}

where $r_0$ is the galaxy scale radius. Since $r_0$ varies between galaxies, 
$a_0$ is not truly universal in RS — it only appears universal because of the 
correlation between $r_0$ and galaxy luminosity.

\section{Laboratory Tests}

\begin{remark}[Laboratory Undetectability]
At laboratory scales ($r \ll r_0$), the ILG correction is
\begin{equation}
\delta g / g = \alpha_t \cdot (r/r_0) \approx 0.191 \times (10~\mathrm{m} / 10~\mathrm{kpc}) 
\approx 6 \times 10^{-22},
\end{equation}
completely undetectable in laboratory tests --- consistent with precision tests of 
Newtonian gravity at $10$--$100~\mu$m scales showing no deviation.
\end{remark}

\section{Conclusion}

ILG provides a zero-parameter derivation of galactic rotation curves:
\begin{enumerate}
\item Coupling $\alpha_t = (1 - \varphi^{-1})/2 \approx 0.191$ from $\varphi$
\item $M/L = \varphi \approx 1.618$ solar units (universal)
\item Flat rotation curves at large $r$ (ILG term dominates)
\item Tully-Fisher relation $v_\infty \propto M_b^{1/4}$ automatic
\item UDG diversity from substrate coherence variation
\item No per-galaxy tuning: all galaxies fit with the same $\alpha_t$
\end{enumerate}

The theory predicts universal $\alpha_t = 0.191$ and no dark matter particles. 
These can be tested by analyzing the SPARC sample and searching for variations in $\alpha_t$.

\begin{thebibliography}{9}
\bibitem{sparc} F. Lelli, S.S. McGaugh, J.M. Schombert, \emph{SPARC: Mass Models for 175 Disk Galaxies with Spitzer Photometry and Accurate Rotation Curves}, AJ 152:157 (2016).
\bibitem{milgrom} M. Milgrom, \emph{A modification of the Newtonian dynamics as a possible alternative to the hidden mass hypothesis}, ApJ 270:365 (1983).
\end{thebibliography}

\end{document}
